\documentclass[10pt]{article}
\pagestyle{plain}

\usepackage{amsmath}
\usepackage{newpxtext,newpxmath}
\usepackage[scaled=1.03]{inconsolata}

\usepackage[top=1in, bottom=1in, left=1.25in, right=1.25in]{geometry}
\usepackage[hidelinks]{hyperref}
\usepackage{graphicx}
\usepackage{listings}
	\lstset{basicstyle = \linespread{1.2}\ttfamily, basewidth = {0.5em,0.5em}, breaklines}
\usepackage{parskip}
\usepackage{setspace}
	\setstretch{1.3}

\usepackage{booktabs}
	\def\sym#1{\ifmmode^{#1}\else\(^{#1}\)\fi}
\usepackage{dcolumn}
	\newcolumntype{.}{D{.}{.}{-1}}

\makeatletter
\let\ftype@table\ftype@figure
\makeatother



\begin{document}



Decision Lab, MIMIC-IV v3.1 \\
Barney Guo, November 26, 2026



\subsection*{A logistic regression for delirium in the ICU}
\bigskip
\subsubsection*{Purpose}

The clinician test on the Decision Lab page currently shows illustrative risks. To replace them, we fit a simple prediction model on MIMIC-IV and choose the eight test cases from its real predictions, including cases where it is wrong. This note sets out the model, how it turns a patient's record into a risk, and how it will be checked.

Numbers marked as illustrative are made up to show the arithmetic. Empty table cells are filled in as queries 13--19 in this folder are run.



\subsubsection*{Part 1: Cohort and outcome}

Each row of the data is one ICU stay (queries 09 and 14).

\begin{itemize}
	\item \emph{Who:} each adult patient's first ICU stay that lasted at least 24 hours.
	\item \emph{When the prediction is made:} at $t_0$, 24 hours after ICU admission. Predictors use only data recorded before $t_0$.
	\item \emph{Outcome:} $y_i = 1$ if any CAM-ICU assessment (item 228332) between $t_0$ and $t_0 + 48$~hours is positive, and $y_i = 0$ if the assessments in that window are all negative. Stays with no valid assessment in the window, and stays already positive before $t_0$, are excluded.
	\item \emph{Training and test years:} the model is fit on 2008--2016 and tested on 2017--2019, which mimics using it on new patients after deployment. 2020--2022 is left out because only 42\% of those stays have any medication or procedure record (query 12).
\end{itemize}

\textbf{Table 1.} Cohort flow. Delirium counts for the two splits come from query 14.
\begin{center}
	\begin{tabular*}{\textwidth}{@{\hskip\tabcolsep\extracolsep{\fill}} l r r}
	\toprule
	Stays											& Number	& Delirium in 48 h \\
	\midrule
	First ICU stays of 24 hours or more					& 51,837	& \\
	\quad Excluded: delirium found in the first 24 hours	& 8,079		& \\
	\quad Excluded: no valid assessment in the next 48 hours	& 20,958	& \\
	Cohort, 2008--2022									& 22,800	& 3,535 (15.5\%) \\
	\quad Training, 2008--2016							& 13,520	& 14.9\% \\
	\quad Test, 2017--2019								& 5,495		& 15.7\% \\
	\bottomrule
	\end{tabular*}
\end{center}



\subsubsection*{Part 2: The model}

\emph{In plain terms.} Logistic regression turns what we know about a patient into a risk between 0\% and 100\%. Each characteristic, such as age or being on a ventilator, adds to or subtracts from a score, and a fixed formula converts the score into a probability. Fitting the model means finding how much each characteristic should count, using patients whose outcomes are already known. It is one of the oldest and most widely used methods for clinical risk scores, and its weights can be read directly, which suits a study about how risk is communicated.

Let $p_i$ be the probability that stay $i$ develops delirium in the 48 hours after $t_0$. Logistic regression models the log odds of $p_i$ as a linear function of the predictors:
\begin{equation*}
	\begin{aligned}
	z_i = \log \frac{p_i}{1 - p_i} &= \beta_0 + \beta_1 \, \text{age}_i + \beta_2 \, \text{female}_i \\
	&\quad + \beta_3 \, \text{emergency}_i + \beta_4 \, \text{planned}_i + \textstyle\sum_{k} \gamma_k \, \text{unit}_{ik} \\
	&\quad + \beta_5 \, \text{dementia}_i + \beta_6 \, \text{vent}_i + \beta_7 \, \text{benzo}_i + \beta_8 \, \text{sedative}_i + \beta_9 \, \text{opioid}_i .
	\end{aligned}
\end{equation*}
Solving for $p_i$ turns the score $z_i$ into a risk between 0 and 1:
\begin{equation*}
	p_i = \frac{1}{1 + e^{-z_i}} .
\end{equation*}

Every predictor except age is an indicator equal to 1 or 0. Each text variable (sex, admission group, ICU type) has one reference category that gets no term. The intercept $\beta_0$ is the log odds when every predictor is zero, which includes an age of zero, so it only sets the level of the scale and has no clinical meaning of its own. The equation is written with male, observation admission, and one ICU type as references for readability. BigQuery chooses the references itself when it fits the model, and Table 4 records which categories it left out.

\textbf{Table 2.} Predictors.
\begin{center}
	\begin{tabular}{@{} l p{0.72\textwidth} @{}}
	\toprule
	Predictor			& Definition \\
	\midrule
	age					& Age in years at ICU admission. MIMIC-IV shifts dates and caps the oldest ages, so ages are approximate. \\
	female				& Sex recorded as female. \\
	emergency, planned	& Admission type grouped as emergency (EW EMER., URGENT, DIRECT EMER.), planned (SURGICAL SAME DAY ADMISSION, ELECTIVE), or observation (all other types). \\
	unit				& First care unit of the ICU stay, one indicator per type other than the reference. Types with fewer than 100 training stays are merged into ``Other'': in a first fit, the rarest types produced extreme coefficients, a sign of separation. \\
	dementia			& ICD-10 F01--F03 or G30, or ICD-9 290, 294.1 or 331.0 (excluding 290.3, 290.11 and 290.41), coded in the same hospitalization. \\
	vent				& Invasive ventilation started before $t_0$. \\
	benzo				& Midazolam, lorazepam or diazepam started before $t_0$. \\
	sedative			& Propofol or dexmedetomidine started before $t_0$. \\
	opioid				& Fentanyl, morphine or hydromorphone started before $t_0$. \\
	\bottomrule
	\end{tabular}
\end{center}

Language, race and insurance are kept for the subgroup checks in Part 5 but are not predictors.



\subsubsection*{Part 3: Fitting the model}

The coefficients are chosen to minimize the log loss on the training stays, which is the same as maximizing the likelihood:
\begin{equation*}
	\hat{\beta} = \arg\min_{\beta} \; -\frac{1}{n} \sum_{i=1}^{n} \bigl[ y_i \log p_i + (1 - y_i) \log (1 - p_i) \bigr] .
\end{equation*}

Query 15 fits the model in BigQuery ML:
\begin{lstlisting}
CREATE OR REPLACE MODEL `physionet-data-513915.delirium.lr_24h`
OPTIONS (
  model_type = 'logistic_reg',
  input_label_cols = ['delirium_48h'],
  data_split_method = 'NO_SPLIT',
  category_encoding_method = 'DUMMY_ENCODING',
  max_iterations = 50,
  min_rel_progress = 0.0001
) AS
SELECT delirium_48h, age, gender, admission_group, first_careunit,
       dementia, ventilated_24h, benzo_24h, sedative_24h, opioid_24h
FROM `physionet-data-513915.delirium.model_data`
WHERE split = 'train';
\end{lstlisting}

\begin{itemize}
	\item \emph{No class weighting.} With about 15\% delirium, weighting the cases would push every predicted risk up, and the test page shows those risks to clinicians.
	\item \emph{No penalty.} L1 and L2 regularization stay at BigQuery's default of zero. With about 2,000 delirium cases in the training years and fewer than 20 coefficients, the fit is stable without one.
	\item \emph{Dummy encoding.} Each text variable loses one category, so each remaining coefficient compares a category with its reference. BigQuery's default one-hot encoding keeps every category, and its coefficients cannot be read that way.
	\item \emph{Convergence.} BigQuery fits by gradient descent and stops once the loss improves by less than 0.01\% in an iteration. It allows at most 50 iterations per run, so query 15 continues training from the saved model with a warm start (\texttt{warm\_start = TRUE}) until a run ends before its cap. Its last statement lists the loss at the final iterations.
	\item \emph{Scale.} BigQuery standardizes age while fitting, and query 19 reports the coefficients on the original scale, per year of age.
\end{itemize}

\textbf{Table 3.} Convergence, from query 15.
\begin{center}
	\begin{tabular*}{0.6 \textwidth}{@{\hskip\tabcolsep\extracolsep{\fill}} l r}
	\toprule
	Iterations run							& 52 (50, then 2 in a warm start) \\
	Training loss at the last iteration		& 0.37224 \\
	Change in the last iteration			& less than 0.00001 \\
	\bottomrule
	\end{tabular*}
\end{center}



\subsubsection*{Part 4: Reading the coefficients}

\begin{itemize}
	\item \emph{Odds ratios.} When an indicator changes from 0 to 1 with everything else fixed, the odds of delirium are multiplied by $e^{\beta_j}$. For age, $e^{\beta_1}$ is per year and $e^{10 \beta_1}$ is per ten years.
	\item \emph{Odds are not risks.} At a baseline risk of 15\%, an odds ratio of 2 raises the risk to 26\%, not 30\%: the odds go from $0.15 / 0.85 = 0.176$ to $0.353$, and $0.353 / 1.353 = 0.261$.
\end{itemize}

\emph{Worked example, with illustrative coefficients.} Suppose $\beta_0 = -3.0$, $\beta_1 = 0.02$, $\beta_6 = 0.8$, and every other term is zero (a man in the reference ICU type, admitted for observation, with no dementia and no sedating drugs). For a 75-year-old ventilated patient,
\begin{equation*}
	z = -3.0 + 0.02 \cdot 75 + 0.8 = -0.7 ,
	\qquad
	p = \frac{1}{1 + e^{0.7}} = 0.332 .
\end{equation*}
The full display on the test page would show this as ``33 in 100 similar patients develop delirium in the next 48 hours.'' The real coefficients come from query 19.

\textbf{Table 4.} Coefficients, from query 19, after merging rare ICU types.
\begin{center}
	\begin{tabular*}{0.8 \textwidth}{@{\hskip\tabcolsep\extracolsep{\fill}} l . .}
	\toprule
	Predictor							& \multicolumn{1}{c}{Log odds, $\hat{\beta}$}	& \multicolumn{1}{c}{Odds ratio, $e^{\hat{\beta}}$} \\
	\midrule
	Intercept	& -4.308 & 0.013 \\
	Age, per year	& 0.019 & 1.02 \\
	Sex: female	& -0.005 & 1.00 \\
	Admission: observation	& 0.040 & 1.04 \\
	Admission: planned	& -0.589 & 0.55 \\
	\midrule
	ICU type: Coronary Care Unit (CCU)	& 0.184 & 1.20 \\
	ICU type: Medical ICU (MICU)	& 1.111 & 3.04 \\
	ICU type: Medical/Surgical ICU (MICU/SICU)	& 0.469 & 1.60 \\
	ICU type: Neuro Intermediate	& 0.852 & 2.34 \\
	ICU type: Neuro Stepdown	& 1.123 & 3.07 \\
	ICU type: Neuro Surgical ICU (Neuro SICU)	& 1.712 & 5.54 \\
	ICU type: Surgical ICU (SICU)	& 0.690 & 1.99 \\
	ICU type: Trauma SICU (TSICU)	& 0.696 & 2.01 \\
	ICU type: Other	& 2.214 & 9.15 \\
	\midrule
	Dementia	& 0.931 & 2.54 \\
	Invasive ventilation	& 1.170 & 3.22 \\
	Benzodiazepine	& 0.660 & 1.94 \\
	Other sedative	& 0.202 & 1.22 \\
	Opioid	& 0.104 & 1.11 \\
	\midrule
	\multicolumn{3}{@{}l}{Reference categories: male, emergency admission, Cardiac Vascular ICU (CVICU).} \\
	\bottomrule
	\end{tabular*}
\end{center}

BigQuery chose emergency admission, not observation, as the admission reference, so the admission rows compare with emergency admissions. The ``Other ICU type merges the types with fewer than 100 training stays, so its coefficient rests on few stays and is imprecise. Odds ratios here are associations with everything else held fixed, not causal effects.

\emph{Worked example, with the fitted coefficients.} For a 75-year-old man admitted as an emergency to the CVICU, with none of the other risk factors, $z = -4.308 + 0.019 \cdot 75 = -2.883$ and $p = 1 / (1 + e^{2.883}) = 0.053$, about 5 in 100. The same man in the MICU on a ventilator who received a benzodiazepine has $z = -2.883 + 1.111 + 1.170 + 0.660 = 0.058$ and $p = 0.514$, about 51 in 100.



\subsubsection*{Part 5: Checking the model}

All checks use the test years, 2017--2019, except the training row of Table 5.

\emph{Discrimination.} In plain terms, the AUROC (area under the receiver operating characteristic curve) measures how well the model puts patients in the right order of risk. Pick one patient who developed delirium and one who did not: the AUROC is how often the model gives the first patient the higher risk. A value of 0.5 is no better than a coin toss and 1.0 is a perfect ordering; values between 0.7 and 0.8 are usually described as acceptable and above 0.8 as excellent. The AUROC says nothing about whether the risk numbers themselves are right, which is what calibration, below, checks. Formally, the AUROC is the probability that a randomly chosen stay with delirium gets a higher predicted risk than a randomly chosen stay without it, with ties counted as half:
\begin{equation*}
	\text{AUROC} = \frac{1}{n_1 n_0} \sum_{i : y_i = 1} \; \sum_{j : y_j = 0} \Bigl( \mathbf{1}[\hat{p}_i > \hat{p}_j] + \tfrac{1}{2} \, \mathbf{1}[\hat{p}_i = \hat{p}_j] \Bigr) ,
\end{equation*}
where $n_1$ and $n_0$ are the numbers of stays with and without delirium. Query 18 computes the same quantity from ranks, $\text{AUROC} = \bigl( R_1 - n_1 (n_1 + 1) / 2 \bigr) / (n_1 n_0)$, where $R_1$ is the sum of the ranks of the stays with delirium and tied stays share their average rank.

\emph{Calibration.} A model is calibrated when its predicted risks match the rates actually observed. Query 17 sorts the test stays into ten equal groups by predicted risk and compares the mean predicted risk with the observed share in each group. The test page shows risks as numbers, so calibration matters as much as the AUROC. In Jauk et al., the AUROC held steady during deployment while calibration drifted.

\textbf{Table 5.} Performance, from query 16.
\begin{center}
	\begin{tabular*}{0.7 \textwidth}{@{\hskip\tabcolsep\extracolsep{\fill}} l r . .}
	\toprule
	Years					& Stays		& \multicolumn{1}{c}{AUROC}	& \multicolumn{1}{c}{Log loss} \\
	\midrule
	Training, 2008--2016	& 13,520	& 0.739 & 0.3723 \\
	Test, 2017--2019		& 5,495		& 0.764 & 0.3723 \\
	\bottomrule
	\end{tabular*}
\end{center}

The training log loss from query 16 (0.3723) differs slightly from the last loss BigQuery reports during fitting (0.37224, Table 3).

\textbf{Table 6.} Calibration in the test years, from query 17. Group 1 has the lowest predicted risks.
\begin{center}
	\begin{tabular*}{0.6 \textwidth}{@{\hskip\tabcolsep\extracolsep{\fill}} c r . .}
	\toprule
	Group	& Stays	& \multicolumn{1}{c}{Predicted (\%)}	& \multicolumn{1}{c}{Observed (\%)} \\
	\midrule
	1		& 550 & 3.8 & 3.3 \\
	2		& 550 & 5.4 & 4.7 \\
	3		& 550 & 6.8 & 6.9 \\
	4		& 550 & 8.3 & 8.0 \\
	5		& 550 & 9.8 & 7.5 \\
	6		& 549 & 11.4 & 14.2 \\
	7		& 549 & 13.5 & 14.8 \\
	8		& 549 & 17.0 & 14.9 \\
	9		& 549 & 23.0 & 31.1 \\
	10		& 549 & 40.7 & 51.5 \\
	\bottomrule
	\end{tabular*}
\end{center}

In the test years, predicted and observed risks agree within about 3 percentage points in groups 1 to 8. The model underestimates the two highest groups: 23.0\% predicted against 31.1\% observed in group 9, and 40.7\% against 51.5\% in group 10. Across all test stays the mean predicted risk is 14.0\% and the observed rate 15.7\%. With about 550 stays per group, the observed rates carry a 95\% interval of roughly $\pm 4$ points in the top two groups, so both gaps are larger than chance alone.

\emph{Subgroups.} Query 18 repeats the checks within groups of patients. A group whose mean predicted risk sits well below its observed rate is underestimated by the model, which makes it a candidate for the test case that carries a warning. AUROCs from fewer than about 30 delirium cases are unstable, so they are reported with their counts.

\textbf{Table 7.} Subgroups in the test years, from query 18. Groups with fewer than 20 stays (language unknown, and the "no charge" insurance category) are not shown.
\begin{center}
	\begin{tabular*}{\textwidth}{@{\hskip\tabcolsep\extracolsep{\fill}} l r r . . .}
	\toprule
	Group					& Stays	& Delirium	& \multicolumn{1}{c}{Observed (\%)}	& \multicolumn{1}{c}{Predicted (\%)}	& \multicolumn{1}{c}{AUROC} \\
	\midrule
	All	& 5,495 & 862 & 15.7 & 14.0 & 0.764 \\
	\midrule
	Age 18--64	& 2,561 & 326 & 12.7 & 11.6 & 0.801 \\
	Age 65--79	& 1,976 & 321 & 16.2 & 14.6 & 0.721 \\
	Age 80+	& 958 & 215 & 22.4 & 18.8 & 0.716 \\
	\midrule
	Dementia	& 205 & 87 & 42.4 & 34.9 & 0.630 \\
	No dementia	& 5,290 & 775 & 14.7 & 13.2 & 0.756 \\
	\midrule
	Sex: F	& 2,320 & 350 & 15.1 & 13.8 & 0.772 \\
	Sex: M	& 3,175 & 512 & 16.1 & 14.1 & 0.759 \\
	\midrule
	Admission: emergency	& 2,929 & 592 & 20.2 & 15.8 & 0.768 \\
	Admission: planned	& 1,000 & 82 & 8.2 & 8.2 & 0.734 \\
	Admission: observation	& 1,566 & 188 & 12.0 & 14.3 & 0.732 \\
	\midrule
	Language: English	& 4,942 & 769 & 15.6 & 13.7 & 0.761 \\
	Language: Other	& 543 & 91 & 16.8 & 16.1 & 0.788 \\
	\midrule
	Race: White	& 3,521 & 507 & 14.4 & 13.3 & 0.741 \\
	Race: Black	& 398 & 62 & 15.6 & 13.8 & 0.768 \\
	Race: Hispanic or Latino	& 174 & 25 & 14.4 & 12.1 & 0.838 \\
	Race: Asian	& 175 & 20 & 11.4 & 13.2 & 0.790 \\
	Race: Other	& 289 & 44 & 15.2 & 12.6 & 0.751 \\
	Race: Unknown	& 938 & 204 & 21.7 & 17.4 & 0.798 \\
	\midrule
	Insurance: Medicare	& 2,740 & 502 & 18.3 & 15.9 & 0.724 \\
	Insurance: Medicaid	& 857 & 138 & 16.1 & 12.7 & 0.798 \\
	Insurance: Private	& 1,676 & 190 & 11.3 & 11.6 & 0.787 \\
	Insurance: Other	& 149 & 20 & 13.4 & 13.5 & 0.850 \\
	Insurance: Unknown	& 72 & 12 & 16.7 & 13.0 & 0.840 \\
	\bottomrule
	\end{tabular*}
\end{center}

Across all test stays the model predicts 14.0\% against 15.7\% observed, a gap of 1.7 points. The largest gap is in patients with coded dementia: 34.9\% predicted against 42.4\% observed, with the lowest AUROC (0.63). Emergency admissions (15.8\% against 20.2\%), patients aged 80 and over (18.8\% against 22.4\%), stays with race recorded as unknown (17.4\% against 21.7\%) and Medicaid patients (12.7\% against 16.1\%) are also underestimated by more than the overall gap. Observation admissions are overestimated (14.3\% against 12.0\%). Patients whose language is not English show little gap (16.1\% against 16.8\%), but a language barrier can block the assessment itself (query 02), so patients left without an outcome are missing from this check. Groups with fewer than about 30 delirium cases give unstable AUROCs.



\subsubsection*{Part 6: Decisions still open}

\begin{itemize}
	\item \emph{Prediction time.} Predict 12 or 24 hours after ICU admission. Table 8 shows when delirium is first detected. Of 14,390 first ICU stays with a positive result, 46.7\% are first positive within 12 hours of admission and 62.8\% within 24 hours, so most delirium in the ICU is already present, or at least first found, before either prediction time. A 24-hour prediction excludes all of these stays. A 12-hour prediction keeps the 2,325 stays first positive at 12--24 hours as possible outcomes, but has only 12 hours of data to predict from. Detection depends on when assessments are done, so some of those 2,325 may be a first assessment of delirium present on arrival. The 24--72-hour window holds 3,544 first positives, close to the 3,535 outcomes in the cohort, which supports how query 09 built it.
	\item \emph{What counts as a high risk.} Choosing test cases where the model is right or wrong needs a cut-off. One option is the top 15\% of predicted risk, which matches the red and yellow icons in Jauk et al.\ (85th percentile and above). Another is a fixed predicted risk agreed with clinicians.
	\item \emph{Outcome window.} The main analysis keeps 48 hours. Sensitivity analyses widen the window and count patients who left the ICU without an assessment as not delirious, which alone lowers the rate from 15.5\% to about 10.2\%.
\end{itemize}

\textbf{Table 8.} When the first positive CAM-ICU appears, among first ICU stays with any positive result, from query 13.
\begin{center}
	\begin{tabular*}{0.6 \textwidth}{@{\hskip\tabcolsep\extracolsep{\fill}} l r .}
	\toprule
	Time after ICU admission	& Stays	& \multicolumn{1}{c}{Share (\%)} \\
	\midrule
	Day 1, 0--12 hours		& 6,713		& 46.7 \\
	Day 1, 12--24 hours		& 2,325		& 16.2 \\
	Day 2, 24--48 hours		& 2,431		& 16.9 \\
	Day 3, 48--72 hours		& 1,113		& 7.7 \\
	Days 4--7				& 1,363		& 9.5 \\
	After day 7				& 445		& 3.1 \\
	\midrule
	All						& 14,390	& 100.0 \\
	\bottomrule
	\end{tabular*}
\end{center}



\subsubsection*{Limitations}

\begin{itemize}
	\item Dementia comes from diagnosis codes for the same hospitalization, which are assigned at discharge. Coded dementia is only 2--4\%, so it is probably under-recorded.
	\item Coding moved from ICD-9 to ICD-10 in October 2015, inside the training years.
	\item CAM-ICU is missing or ``unable to assess'' more often for patients with a language barrier or advanced dementia, so the outcome itself is less complete for them.
	\item All stays come from one academic medical center in Boston.
\end{itemize}



\end{document}
